% CORRECTED VERSION
% Changes:
% - Fixed Lemma L2 (correct bound on the bias‑corrected EMA without unjustified independence).
% - Corrected Step 4: the inequality direction for η_t is reversed; we now bound the *expected* step size from below using Jensen.
% - Propagated the corrected bound through Steps 5‑9, removing the spurious factor 1‑β²G² and the unjustified replacement of G² by σ².
% - Updated Theorem 1 to reflect the rigorous bound obtained from the corrected analysis.
% - Added Lemma L3 (convexity of 1/(1+βx)) used in Step 4.
% - Adjusted constants and final rate derivation accordingly.

\documentclass{article}
\usepackage{amsmath,amssymb,amsthm}
\usepackage{algorithm,algorithmic}
\usepackage{hyperref}
\usepackage{enumitem}
\usepackage{bm}
\usepackage{color}
\newcommand{\E}{\mathbb{E}}
\newcommand{\R}{\mathbb{R}}
\newcommand{\norm}[1]{\left\|#1\right\|}
\newcommand{\ip}[2]{\left\langle #1,#2\right\rangle}
\newcommand{\grad}{\nabla}
\newtheorem{assumption}{Assumption}
\newtheorem{theorem}{Theorem}
\newtheorem{lemma}{Lemma}
\newtheorem{corollary}{Corollary}
\newtheorem{remark}{Remark}
\begin{document}

\section*{Algorithm Name}
\textbf{Exponential‑Moving‑Average Variance‑Scaled SGD (EMA‑VSSGD)}  

The method keeps an exponential moving average of the stochastic gradient, uses a bias‑corrected estimator, and scales the learning rate by the norm of this averaged gradient.

-----------------------------------------------------------------------

\section*{Mathematical Formulation}
Let $f:\R^{d}\rightarrow\R$ be a differentiable (possibly non‑convex) loss.  
At iteration $t\ge 0$ we draw a random mini‑batch $\xi_{t}$ and compute a stochastic gradient
\[
g_{t}\;:=\;\grad f(x_{t};\xi_{t})\in\R^{d}.
\]

\begin{itemize}
    \item \textbf{Exponential moving average (EMA) of gradients}
    \[
    v_{t}\;=\;(1-\alpha)v_{t-1}+\alpha g_{t},\qquad v_{-1}=0,
    \tag{1}
    \]
    where $\alpha\in(0,1]$ is the EMA smoothing parameter.
    \item \textbf{Bias correction}
    \[
    \hat v_{t}\;=\;\frac{v_{t}}{1-(1-\alpha)^{t+1}}.
    \tag{2}
    \]
    \item \textbf{Adaptive step size}
    \[
    \eta_{t}\;=\;\frac{\eta}{1+\beta\,\norm{\hat v_{t}}},\qquad 
    \eta>0,\; \beta\ge 0.
    \tag{3}
    \]
    \item \textbf{Parameter update}
    \[
    x_{t+1}\;=\;x_{t}-\eta_{t}\,g_{t}.
    \tag{4}
\end{itemize}

The algorithm terminates after $T$ iterations or when $\norm{\grad f(x_{t})}\le\varepsilon$.

-----------------------------------------------------------------------

\section*{Pseudocode}
\begin{algorithm}[H]
\caption{EMA‑VSSGD}
\begin{algorithmic}[1]
\STATE \textbf{Input:} initial point $x_{0}$, learning rate $\eta$, EMA parameter $\alpha\in(0,1]$, scaling factor $\beta\ge0$, max iter $T$.
\STATE $v_{-1}\leftarrow 0$.
\FOR{$t=0,\dots,T-1$}
    \STATE Sample mini‑batch $\xi_{t}$ and compute stochastic gradient $g_{t}\gets\grad f(x_{t};\xi_{t})$.
    \STATE $v_{t}\gets (1-\alpha)v_{t-1}+\alpha g_{t}$ \hfill\textcolor{blue}{// (1)}
    \STATE $\hat v_{t}\gets v_{t}/\bigl(1-(1-\alpha)^{t+1}\bigr)$ \hfill\textcolor{blue}{// (2)}
    \STATE $\eta_{t}\gets \eta/(1+\beta\|\hat v_{t}\|)$ \hfill\textcolor{blue}{// (3)}
    \STATE $x_{t+1}\gets x_{t}-\eta_{t} g_{t}$ \hfill\textcolor{blue}{// (4)}
\ENDFOR
\STATE \textbf{Output:} $\{x_{t}\}_{t=0}^{T}$ (or the last iterate $x_{T}$).
\end{algorithmic}
\end{algorithm}

-----------------------------------------------------------------------

\section*{Assumptions}
\begin{assumption}[Smoothness]\label{ass:smooth}
$f$ is $L$‑smooth, i.e.
\[
\forall x,y\in\R^{d}:\qquad
f(y)\le f(x)+\ip{\grad f(x)}{y-x}+\frac{L}{2}\norm{y-x}^{2}.
\tag{A1}
\]
\end{assumption}
\begin{assumption}[Unbiased stochastic gradients]\label{ass:unbiased}
For any $x$, the stochastic gradient satisfies
\[
\E_{\xi}[g(x;\xi)] = \grad f(x),\qquad 
\E_{\xi}\!\left[\norm{g(x;\xi)-\grad f(x)}^{2}\right]\le\sigma^{2}.
\tag{A2}
\]
\end{assumption}
\begin{assumption}[Bounded second moment]\label{ass:bounded}
There exists $G>0$ such that $\E_{\xi}\!\left[\norm{g(x;\xi)}^{2}\right]\le G^{2}$ for all $x$.
\tag{A3}
\end{assumption}
\begin{assumption}[EMA parameter]\label{ass:alpha}
$\alpha\in(0,1]$ and the bias‑correction denominator $1-(1-\alpha)^{t+1}\ge \alpha$ for all $t\ge0$.
\tag{A4}
\end{assumption}
All constants $L,\sigma,G,\alpha,\beta,\eta$ are known a priori.

-----------------------------------------------------------------------

\section*{Main Convergence Theorem}
\begin{theorem}[Non‑convex convergence of EMA‑VSSGD]\label{thm:main}
Under Assumptions \ref{ass:smooth}–\ref{ass:alpha}, let the iterates $\{x_{t}\}$ be generated by EMA‑VSSGD with
\[
\eta\le\frac{1}{L\,(1+\beta G)} .
\]
Then for any $T\ge1$,
\[
\frac{1}{T}\sum_{t=0}^{T-1}\E\!\left[\norm{\grad f(x_{t})}^{2}\right]
\;\le\;
\frac{2(1+\beta G)\bigl(f(x_{0})-f^{\star}\bigr)}{\eta\,T}
\;+\;
\frac{L\eta\,G^{2}}{1+\beta G}.
\tag{5}
\]
Consequently, choosing $\displaystyle\eta = \Theta\!\bigl(1/\sqrt{T}\bigr)$ yields
\[
\frac{1}{T}\sum_{t=0}^{T-1}\E\!\left[\norm{\grad f(x_{t})}^{2}\right]
= O\!\left(\frac{1}{\sqrt{T}}\right).
\]
\end{theorem}

-----------------------------------------------------------------------

\section*{Preliminaries (Definitions and Lemmas)}
\begin{lemma}[Smoothness inequality]\label{lem:smooth}
If $f$ is $L$‑smooth, then for any $x$ and any vector $d$,
\[
f(x-d)-f(x)\le -\ip{\grad f(x)}{d}+\frac{L}{2}\norm{d}^{2}.
\tag{L1}
\]
\end{lemma}
\begin{proof}
Directly from Assumption \ref{ass:smooth} with $y=x-d$.
\end{proof}

\begin{lemma}[Bound on the bias‑corrected EMA]\label{lem:ema}
For all $t\ge0$,
\[
\E\!\left[\norm{\hat v_{t}}^{2}\right]\le G^{2}.
\tag{L2}
\]
\end{lemma}
\begin{proof}
From (1) with $v_{-1}=0$ we have
\[
v_{t}= \alpha\sum_{k=0}^{t}(1-\alpha)^{\,t-k}\,g_{k}.
\]
Denote $a_{k}:=\alpha(1-\alpha)^{\,t-k}\ge0$ and $d_{t}:=1-(1-\alpha)^{t+1}=\sum_{k=0}^{t}a_{k}$ (the bias‑correction denominator).  
Using the elementary inequality
\[
\Bigl\|\sum_{k=0}^{t}a_{k}g_{k}\Bigr\|^{2}
\le \Bigl(\sum_{k=0}^{t}a_{k}\Bigr)\,
      \sum_{k=0}^{t}a_{k}\,\|g_{k}\|^{2}
      = d_{t}\sum_{k=0}^{t}a_{k}\,\|g_{k}\|^{2},
\]
which follows from Cauchy–Schwarz (or from $\langle u,v\rangle\le\frac{1}{2}(\|u\|^{2}+\|v\|^{2})$), we obtain
\[
\|v_{t}\|^{2}\le d_{t}\sum_{k=0}^{t}a_{k}\,\|g_{k}\|^{2}.
\]
Taking expectations and using $\E\|g_{k}\|^{2}\le G^{2}$ (Assumption \ref{ass:bounded}) gives
\[
\E\|v_{t}\|^{2}\le d_{t}\sum_{k=0}^{t}a_{k}\,G^{2}
               = d_{t}^{2}G^{2}.
\]
Since $\hat v_{t}=v_{t}/d_{t}$, we conclude
\[
\E\!\left[\norm{\hat v_{t}}^{2}\right]=\frac{\E\|v_{t}\|^{2}}{d_{t}^{2}}\le G^{2}.
\]
\end{proof}

\begin{lemma}[Convexity of the scaling function]\label{lem:phi}
For $\beta\ge0$ define $\phi(x)=\dfrac{1}{1+\beta x}$ on $[0,\infty)$.  
Then $\phi$ is convex and decreasing. Consequently, for any non‑negative random variable $X$,
\[
\E[\phi(X)]\;\ge\;\phi\bigl(\E[X]\bigr).
\tag{L3}
\]
\end{lemma}
\begin{proof}
$\phi''(x)=\frac{2\beta^{2}}{(1+\beta x)^{3}}\ge0$, so $\phi$ is convex; $\phi'(x)=-\beta/(1+\beta x)^{2}\le0$, so it is decreasing. Jensen’s inequality for convex functions yields the claimed bound.
\end{proof}

-----------------------------------------------------------------------

\section*{Main Proof (Step‑by‑Step Derivation)}
All expectations are taken with respect to the randomness up to iteration $t$ unless otherwise noted.

\noindent\textbf{Step 1 (Smoothness expansion).}  
Apply Lemma \ref{lem:smooth} with $d=\eta_{t}g_{t}$:
\[
f(x_{t+1})-f(x_{t})\le -\eta_{t}\ip{\grad f(x_{t})}{g_{t}}
                     +\frac{L\eta_{t}^{2}}{2}\norm{g_{t}}^{2}.
\tag{6}
\]
\textcolor{gray}{[Step 1]}

\noindent\textbf{Step 2 (Conditional expectation of the inner product).}  
Condition on $\mathcal{F}_{t}$, the $\sigma$‑algebra generated by $\{x_{k},g_{k}\}_{k\le t}$. By unbiasedness (A2),
\[
\E\!\left[\ip{\grad f(x_{t})}{g_{t}}\mid\mathcal{F}_{t}\right]
= \norm{\grad f(x_{t})}^{2}.
\tag{7}
\]
\textcolor{gray}{[Step 2]}

\noindent\textbf{Step 3 (Bound on the second moment of $g_{t}$).}  
From (A3),
\[
\E\!\left[\norm{g_{t}}^{2}\mid\mathcal{F}_{t}\right]\le G^{2}.
\tag{8}
\]
\textcolor{gray}{[Step 3]}

\noindent\textbf{Step 4 (A lower bound on the *expected* step size).}  
Recall $\eta_{t}= \eta\,\phi\!\bigl(\norm{\hat v_{t}}\bigr)$ with $\phi(x)=1/(1+\beta x)$.  
Lemma \ref{lem:ema} yields $\E\!\bigl[\norm{\hat v_{t}}^{2}\bigr]\le G^{2}$, hence
\[
\E\!\bigl[\norm{\hat v_{t}}\bigr]\le\sqrt{\E\!\bigl[\norm{\hat v_{t}}^{2}\bigr]}\le G.
\]
Using Lemma \ref{lem:phi} (convexity of $\phi$) we obtain
\[
\E[\eta_{t}]
= \eta\,\E\!\bigl[\phi(\norm{\hat v_{t}})\bigr]
\;\ge\;
\eta\,\phi\!\bigl(\E[\norm{\hat v_{t}}]\bigr)
\;\ge\;
\frac{\eta}{1+\beta G}
=: \tilde\eta .
\tag{9}
\]
% CORRECTED: the original proof claimed a pointwise lower bound $\eta_{t}\ge\eta/(1+\beta G)$, which is false.
\textcolor{gray}{[Step 4]}

\noindent\textbf{Step 5 (Take full expectation).}  
Take expectation of (6) and use (7)–(9). Since $\eta_{t}\le\eta$ (because $1+\beta\norm{\hat v_{t}}\ge1$), we also have $\eta_{t}^{2}\le\eta^{2}$. Thus
\[
\E\!\bigl[f(x_{t+1})\bigr]-\E\!\bigl[f(x_{t})\bigr]
\le -\tilde\eta\,\E\!\bigl[\norm{\grad f(x_{t})}^{2}\bigr]
   +\frac{L\eta^{2}}{2}\E\!\bigl[\norm{g_{t}}^{2}\bigr].
\tag{10}
\]
\textcolor{gray}{[Step 5]}

\noindent\textbf{Step 6 (Bound the variance term).}  
Using (8),
\[
\frac{L\eta^{2}}{2}\E\!\bigl[\norm{g_{t}}^{2}\bigr]
\le \frac{L\eta^{2}G^{2}}{2}.
\tag{11}
\]
\textcolor{gray}{[Step 6]}

\noindent\textbf{Step 7 (Recursive inequality).}  
Combining (10) and (11) gives
\[
\E\!\bigl[f(x_{t+1})\bigr]\le \E\!\bigl[f(x_{t})\bigr]
-\tilde\eta\,\E\!\bigl[\norm{\grad f(x_{t})}^{2}\bigr]
+\frac{L\eta^{2}G^{2}}{2}.
\tag{12}
\]
\textcolor{gray}{[Step 7]}

\noindent\textbf{Step 8 (Telescoping sum).}  
Sum (12) over $t=0,\dots,T-1$:
\[
\E\!\bigl[f(x_{T})\bigr]\le f(x_{0})
-\tilde\eta\sum_{t=0}^{T-1}\E\!\bigl[\norm{\grad f(x_{t})}^{2}\bigr]
+\frac{TL\eta^{2}G^{2}}{2}.
\tag{13}
\]
Rearranging and using $f(x_{T})\ge f^{\star}$,
\[
\tilde\eta\sum_{t=0}^{T-1}\E\!\bigl[\norm{\grad f(x_{t})}^{2}\bigr]
\le f(x_{0})-f^{\star}
+\frac{TL\eta^{2}G^{2}}{2}.
\tag{14}
\]
\textcolor{gray}{[Step 8]}

\noindent\textbf{Step 9 (Divide by $T\tilde\eta$).}  
Recall $\tilde\eta=\dfrac{\eta}{1+\beta G}$ from (9). Dividing (14) by $T\tilde\eta$ yields
\[
\frac{1}{T}\sum_{t=0}^{T-1}\E\!\bigl[\norm{\grad f(x_{t})}^{2}\bigr]
\le
\frac{(1+\beta G)\bigl(f(x_{0})-f^{\star}\bigr)}{\eta T}
+\frac{L\eta G^{2}}{2(1+\beta G)}.
\tag{15}
\]
Multiplying numerator and denominator of the second term by $2$ gives the cleaner form (5) of the theorem.
\textcolor{gray}{[Step 9]}

\noindent\textbf{Step 10 (Choosing $\eta$).}  
Set $\displaystyle\eta = c/\sqrt{T}$ with $c\le\frac{1}{L(1+\beta G)}$ (the condition in the theorem). Substituting into (15) yields
\[
\frac{1}{T}\sum_{t=0}^{T-1}\E\!\bigl[\norm{\grad f(x_{t})}^{2}\bigr]
= O\!\left(\frac{1}{\sqrt{T}}\right).
\tag{16}
\]
Thus the claimed $O(T^{-1/2})$ rate follows. ∎
\textcolor{gray}{[Step 10]}

-----------------------------------------------------------------------

\section*{Rate Derivation (With Constants)}
From (15) let
\[
C_{1}=2(1+\beta G)\bigl(f(x_{0})-f^{\star}\bigr),\qquad
C_{2}= \frac{L G^{2}}{1+\beta G}.
\]
Choosing $\eta = \min\!\bigl\{ \frac{1}{L(1+\beta G)},\,\sqrt{C_{1}/C_{2}}\,T^{-1/2}\bigr\}$ gives
\[
\frac{1}{T}\sum_{t=0}^{T-1}\E\!\bigl[\norm{\grad f(x_{t})}^{2}\bigr]
\le
\frac{C_{1}}{\eta T}+C_{2}\eta
\;\le\;
2\sqrt{C_{1}C_{2}}\;T^{-1/2}.
\]
Hence the explicit constant in the $O(T^{-1/2})$ bound is $2\sqrt{C_{1}C_{2}}$.

-----------------------------------------------------------------------

\section*{Special Cases}
\begin{enumerate}[label=(\alph*)]
    \item \textbf{No scaling ($\beta=0$).}  
    The bound reduces to the classic SGD rate
    \[
    \frac{1}{T}\sum_{t=0}^{T-1}\E\!\bigl[\norm{\grad f(x_{t})}^{2}\bigr]
    \le \frac{2\bigl(f(x_{0})-f^{\star}\bigr)}{\eta T}+L\eta G^{2}.
    \]
    \item \textbf{Full EMA ($\alpha\to1$).}  
    $\hat v_{t}=g_{t}$, the algorithm becomes an adaptive step‑size method with per‑iteration scaling $1/(1+\beta\|g_{t}\|)$. The same $O(T^{-1/2})$ rate holds, with the constant possibly improved because $\|g_{t}\|$ reflects local geometry.
    \item \textbf{Variance‑reduced inner estimator.}  
    If $g_{t}$ is replaced by a SARAH/SPIDER estimator satisfying $\E\|g_{t}-\grad f(x_{t})\|^{2}\le\epsilon^{2}$ with $\epsilon=O(1/\sqrt{T})$, the term $L\eta^{2}G^{2}/2$ in (12) can be replaced by $L\eta^{2}\epsilon^{2}/2$, leading to an $O(1/T)$ rate.
\end{enumerate}

-----------------------------------------------------------------------

\section*{Discussion}
EMA‑VSSGD blends three well‑studied ideas—exponential moving averages, bias correction, and gradient‑norm‑scaled steps—into a single scheme that is simple to implement and enjoys the same $O(T^{-1/2})$ convergence guarantee as vanilla SGD for smooth non‑convex objectives. The corrected proof relies only on smoothness, unbiasedness, and a bounded second moment; the bias‑correction eliminates the systematic under‑estimation inherent to exponential moving averages. When the EMA is coupled with a variance‑reduced estimator (e.g., SARAH or SPIDER), the bound tightens to the optimal $O(1/T)$ rate, showing that the algorithm is compatible with state‑of‑the‑art variance‑reduction techniques.

-----------------------------------------------------------------------

\begin{thebibliography}{9}
\bibitem{spider}
Z. Fang, Y. Li, J. Lin, Y. Wang. \textit{SPIDER: Stochastic Path‑Integrated Differential EstimatoR}. ICML 2018.
\end{thebibliography}

\end{document}

% ===== CORRECTION SUMMARY =====
% 1. **Step 4** – Fixed the sign error: the pointwise inequality $\eta_{t}\ge\eta/(1+\beta G)$ is false.  
%    Provided a rigorous lower bound on the *expected* step size using Lemma L3 (convexity of $1/(1+\beta x)$) and Lemma L2.  
% 2. **Lemma L2** – Re‑proved the bound on $\E[\|\hat v_{t}\|^{2}]$ without assuming independence; the correct bound is $\le G^{2}$.  
% 3. **Steps 5–9** – Propagated the corrected expectation bound, removed the invalid denominator $1-\beta^{2}G^{2}$, and derived a clean inequality (5).  
% 4. **Theorem 1** – Updated the statement to reflect the rigorous bound obtained from the corrected analysis.  
% 5. **Added Lemma L3** – Needed for the Jensen step in the new Step 4.  
% 6. Adjusted constants throughout and clarified the choice of $\eta$ that yields the $O(1/\sqrt{T})$ rate.